\chapter{Data Residency and On-Chip Memory Design}
\label{chap:ch06}
\typeout{START_CHAPTER "ch06" \theabspage}

\section{All Decode Optimization Is Residency Optimization}
\label{sec:industrial_hook}

A team fused twelve decoder ops on H100 and celebrated lower launches/token---then bytes/token barely moved because every intermediate still exported to HBM between fused boundaries \citep{taxbreak2026,memoryfloor2026}.
Another team fit a layer in Hopper shared memory, shipped the same IR to XDNA, and watched compile fail---identical math, illegal peak footprint \citep{yirage2026,amd2024xdna,nvidia2022hopper}.

\textbf{Every decode bottleneck in this book eventually reduces to residency:} which bytes stay on chip, for how long, and at what tier \citep{dao2022flashattention,chen2020compilerbasedhardw}.
Parts~I--II diagnosed bandwidth and orchestration; Chapters~3--5 named per-class hardware mechanisms \citep{patel2025llminference}.
\textbf{This chapter is the design methodology hub}---how to plan lifetimes, compute peak footprint, and adapt the same fused layer to GPU, CPU, and NPU without repeating hardware catalogs \citep{yirage2026,dlbook}.

\paragraph{Bridge from Chapters~3--5.}
GPU tier sizes, TMA, and bank conflicts: see Chapter~\ref{chap:ch04} \citep{nvidia2022hopper,semianalysis2024tma}.
XDNA static tiles and DMA chains: see Chapter~\ref{chap:ch05} \citep{amd2024xdna,amd2024aie}.
Constraint matrix rows: see Chapter~\ref{chap:ch03} \citep{chen2020compilerbasedhardw}.
Here we only use \emph{capacity numbers} as budget inputs---not re-derive silicon \citep{yirage2026}.

\paragraph{What you will be able to do.}
After this chapter you can (1)~run the four-step residency workflow, (2)~compute peak footprint with double-buffer accounting, (3)~fill a tri-hardware allocation table for one layer, (4)~interpret YiRage memory-pass failures, and (5)~review merges with the residency checklist \citep{yirage2026,taxbreak2026,memoryfloor2026}.

\paragraph{Organization map.}
Lifetime method (\S\ref{sec:ch06_lifetime_planning}), tri-hardware adaptation (\S\ref{sec:ch06_tri_hardware_residency}), worked example (\S\ref{sec:ch06_cross_hw_worked_example}), YiRage pass (\S\ref{sec:ch06_yirage_memory_pass}), engineering patterns (\S\ref{sec:ch06_residency_patterns}), review workflow (\S\ref{sec:ch06_residency_review}), pitfalls (\S\ref{sec:ch06_residency_pitfalls}), takeaways and bridge to Chapter~\ref{chap:ch07} \citep{patel2025llminference}.

\paragraph{Reader paths.}
\textbf{Kernel engineers} should read \S\ref{sec:ch06_lifetime_planning}, the worked example, and pitfalls first---then skim YiRage pass I/O \citep{yirage2026,dao2022flashattention}.
\textbf{Compiler authors} should read the memory pass and pattern sections end-to-end \citep{wu2024mirage,rotem2018glow,chen2020compilerbasedhardw}.
\textbf{Serving leads} should read the industrial hook, tri-hardware table, and bytes/token Pattern~6 \citep{patel2025llminference,memoryfloor2026,taxbreak2026}.
All paths converge on the checklist before Part~IV \citep{mpk2025,yirage2026}.

\paragraph{Industrial narrative.}
TaxBreak shows hundreds of launches per token when fusion boundaries export activations \citep{taxbreak2026}.
Memory-Floor shows bytes/token dominating batch-1 decode even when FLOPs counters look healthy \citep{memoryfloor2026}.
Residency planning attacks the bytes side of that ledger---not peak TFLOPs \citep{dao2022flashattention,patel2025llminference}.

\paragraph{Anti-pattern preview.}
Section~\ref{sec:ch06_residency_pitfalls} collects five failure modes---use them in design review before MegaKernel chapters \citep{yirage2026,taxbreak2026}.
Culture: no fusion merge without a peak-footprint worksheet \citep{chen2020compilerbasedhardw,memoryfloor2026}.
The most expensive mistake is celebrating launch-count drops while bytes/token stays flat---that is residency failure wearing a fusion costume \citep{taxbreak2026,patel2025llminference,memoryfloor2026}.

\paragraph{Metrics reminder.}
Chapter~1 defined launches/token and bytes/token as complementary diagnostics \citep{taxbreak2026,memoryfloor2026}.
This chapter attacks bytes/token by keeping activations and hot weights in fast tiers across fused regions \citep{dao2022flashattention,patel2025llminference}.
When you finish the worksheets here, re-run the same cells as Chapters~4--5 and confirm bytes moved---not just ms/token moved \citep{yirage2026,nvidia2022hopper,amd2024aie}.

\paragraph{Who should read this chapter.}
Compiler authors, kernel engineers porting across GPU/CPU/NPU, and serving teams running triplet fleets \citep{yirage2026,patel2025llminference,amd2024xdna,nvidia2022hopper}.
If you only ship CUDA, still read the tri-hardware table---it explains why your schedule fails elsewhere \citep{chen2020compilerbasedhardw}.

\paragraph{What this chapter is not.}
Not a repeat of Chapter~4 Hopper specs or Chapter~5 XDNA tile catalogs \citep{nvidia2022hopper,amd2024xdna}.
We teach \emph{how to plan} residency given caps already documented upstream \citep{yirage2026,dlbook}.

\paragraph{Why residency is the decode lever.}
Parts~I--II established that batch-1 decode is bandwidth- and orchestration-bound, not FLOP-bound \citep{patel2025llminference,memoryfloor2026}.
TaxBreak-style counts show hundreds to thousands of kernel launches per token when fusion boundaries export activations; each export is a tensor that left the chip and came back \citep{taxbreak2026}.
Removing host overhead alone bought measurable speedups on fused decode cells, but it could not fix illegal HBM traffic---only residency design can \citep{memoryfloor2026,nvidia2024cudagraphs}.
The numbers motivate this chapter; the \emph{method} is what you carry to every target \citep{yirage2026,dao2022flashattention}.

\paragraph{Decode makes lifetime planning tractable.}
At batch-1 the per-token working set is small and regular: a query vector, weight tiles for the current layer, and the KV slice for the current head \citep{kwon2023vllm}.
That regularity means you can solve residency statically at compile time instead of discovering overflow at runtime \citep{rotem2018glow,yirage2026}.
It also means mistakes are expensive: any byte that lives one tier slower than necessary pays its penalty on every subsequent token in the sequence \citep{patel2025llminference,dao2022flashattention}.

\paragraph{Fleet triplet mindset.}
Production fleets rarely ship a single silicon class \citep{patel2025llminference,yirage2026}.
YiRage's value proposition is one fused IR with per-backend residency metadata---not one CUDA schedule pasted onto XDNA \citep{amd2024xdna,chen2020compilerbasedhardw}.
Section~\ref{sec:ch06_cross_hw_worked_example} is the template: same layer, three legal tier assignments, three peak worksheets \citep{yirage2026,memoryfloor2026}.

\index{Data residency}
\index{On-chip memory}
\index{Lifetime planning}

\section{Lifetime Planning: The Four-Step Residency Workflow}
\label{sec:ch06_lifetime_planning}

On-chip residency is the kernel's contract with hardware---planned before codegen, per target \citep{dlbook,yirage2026}.
The mistake is treating fast memory as a post-hoc tweak after the kernel ``works'' \citep{taxbreak2026,memoryfloor2026}.
The first discipline is \textbf{lifetime} planning: for every intermediate tensor, determine when it is produced, when it is last consumed, and therefore where it should live for that span \citep{rotem2018glow,wu2024mirage}.
A tensor whose producer and consumer are adjacent in a fused region never needs device memory; a tensor consumed many operators later has a longer lifetime and may be forced to a slower tier \citep{dao2022flashattention,dlbook}.
Getting this wrong is the most common cause of ``mysterious'' decode regressions---a tensor with an unnecessarily long lifetime occupies fast memory another tensor needed, or spills to HBM and reintroduces the bandwidth tax fusion was meant to remove \citep{taxbreak2026,memoryfloor2026}.

\paragraph{Four-step workflow.}
\begin{enumerate}
\item \textbf{Tensor lifetime inventory}---for each intermediate, mark production point and last consumer on the fused DAG \citep{rotem2018glow,wu2024mirage}.
\item \textbf{Reuse-per-byte ranking}---score tensors by (reuse count)/(bytes); higher score wins faster tiers \citep{dao2022flashattention,dlbook}.
\item \textbf{Peak footprint calculation}---at each timestep, sum simultaneously live buffers; include \emph{double} size for ping-pong operands \citep{semianalysis2024tma,dao2022flashattention}.
\item \textbf{Tier assignment and spill policy}---place from fastest tier down; overflow triggers fusion split or demotion with documented bandwidth cost \citep{yirage2026,chen2020compilerbasedhardw}.
\end{enumerate}

% AUTO_VISUAL:fig_lifetime_planning_pipeline
\begin{figure}[t]
\centering
\begin{tikzpicture}[vis base, scale=0.75, transform shape]
 \node[vis pipeline node, minimum width=2.4cm] (s1) at (0,2.2) {\footnotesize 1. Lifetime\\inventory};
 \node[vis arch node, minimum width=2.4cm] (s2) at (0,1.1) {\footnotesize 2. Reuse\\ranking};
 \node[vis arch node, minimum width=2.4cm] (s3) at (0,0.0) {\footnotesize 3. Peak\\footprint};
 \node[vis pipeline node, minimum width=2.4cm] (s4) at (0,-1.1) {\footnotesize 4. Tier\\assign};
 \draw[vis arrow] (s1) -- (s2) -- (s3) -- (s4);
\end{tikzpicture}
\caption{Residency design four-step workflow (compile-time).}
\label{fig:ch06_residency_workflow}
\end{figure}

\begin{figure}[t]
\centering
\begin{tikzpicture}[vis base, scale=0.88, transform shape]
 \node[vis pipeline node] (s0) at (-5.03,0) {\footnotesize Tensor\\produced};
 \node[vis arch node] (s1) at (-1.68,0) {\footnotesize Assign\\tier};
 \node[vis pipeline node] (s2) at (1.67,0) {\footnotesize Consumed\\on-chip};
 \node[vis arch node] (s3) at (5.03,0) {\footnotesize Buffer\\reused};
 \draw[vis arrow] (s0.east) -- (s1.west);
 \draw[vis arrow] (s1.east) -- (s2.west);
 \draw[vis arrow] (s2.east) -- (s3.west);
\end{tikzpicture}
\caption{Lifetime span: production to last use defines minimum residency duration.}
\label{fig:lifetime_planning_pipeline}
\end{figure}

\paragraph{KV cache lifetime (special case).}
Unlike short-lived activations, KV is written once and read every subsequent token---lifetime spans the whole generation \citep{kwon2023vllm,dao2023flashdecoding}.
Design question: how much KV streams per step and how overlap hides load latency---not whether KV ``fits'' in the fastest tier \citep{patel2025llminference}.
A residency plan that treats the KV slice like a short-lived activation will either exhaust fast memory or stall on loads \citep{dao2022flashattention,memoryfloor2026}.
Paged KV in serving stacks (Chapter~1) decides DRAM placement; this section decides how much KV enters the fast tier each decode step \citep{kwon2023vllm,yirage2026}.

\paragraph{Double-buffer decision rule.}
Enable ping-pong when streamed operand latency exceeds compute on that stage \emph{and} doubled footprint still fits tier budget \citep{semianalysis2024tma,dao2022flashattention}.
The lifetime planner must account for doubled footprint when checking capacity---the first place naive plans diverge from real ones \citep{yirage2026,memoryfloor2026}.
If step~3 peak exceeds capacity only after doubling, shrink tile or split fusion---do not ``enable async copy and hope'' \citep{semianalysis2024tma,taxbreak2026}.

\paragraph{Schedule shapes lifetimes.}
The same dataflow graph can be scheduled with short lifetimes (producer adjacent to consumer) or stretched lifetimes (unrelated work interleaved) \citep{wu2024mirage,dlbook}.
Residency planning and scheduling co-optimize: minimize launches without stretching critical tensors past fast-tier capacity \citep{taxbreak2026,memoryfloor2026}.
The planner chooses an execution order that keeps highest-reuse tensors' lifetimes short enough to stay resident while exposing overlap to hide streaming latency \citep{patel2025llminference,dao2022flashattention}.

\paragraph{Static allocation payoff.}
Known lifetimes enable buffer reuse across non-overlapping intervals---no per-token allocator in the hot path \citep{rotem2018glow,yirage2026}.
Peak footprint is provable before run; compare against Chapter~3 chip-model caps \citep{chen2020compilerbasedhardw}.
Megakernels depend on this property: if peak is unknown, the fused region cannot ship \citep{wu2024mirage,mpk2025,taxbreak2026}.

\paragraph{Reuse-per-byte scoring (step 2 detail).}
Score $s = (\text{reuse count}) / (\text{bytes})$ for each tensor \citep{dlbook,dao2022flashattention}.
Softmax carriers and QKV fragments typically score highest in decode layers---assign them to the fastest tier that fits their \emph{combined} peak interval \citep{dao2022flashattention,patel2025llminference}.
Low-score temporaries demote first when step~4 overflows \citep{yirage2026,wu2024mirage}.

\paragraph{Peak footprint walkthrough (step 3 detail).}
Walk the fused region in program order; maintain a live set; add buffer bytes on allocate, subtract on last use \citep{rotem2018glow,yirage2026}.
For each streamed operand with ping-pong, add $2 \times$ tile bytes to the live set during the overlap window \citep{semianalysis2024tma,dao2022flashattention}.
The maximum live-set size over the walk is the peak footprint---compare to tier cap \citep{nvidia2022hopper,amd2024aie}.

\paragraph{Tier assignment policy (step 4 detail).}
Fill fastest tier with highest-score tensors until full \citep{dlbook}.
On overflow: (a)~split fusion at lowest-score export edge, or (b)~demote lowest-score live tensor to next tier and record expected bytes/token cost \citep{chen2020compilerbasedhardw,yirage2026,memoryfloor2026}.
Never demote carriers or QKV tiles before demoting optional scratch \citep{dao2022flashattention,patel2025llminference}.
Document every demotion in IR metadata so later passes and humans see the bandwidth trade explicitly \citep{yirage2026,wu2024mirage}.

\paragraph{Step 1 inventory (how-to).}
Start from the fused subgraph IR: walk ops in topological order \citep{rotem2018glow,chen2020compilerbasedhardw}.
For each SSA value that materializes bytes, record: name, dtype, shape product, defining op, and uses---last use defines lifetime end \citep{dlbook,yirage2026}.
Export edges (Chapter~2) are lifetime extenders: any tensor written to global mem before its last on-chip consumer effectively has infinite lifetime from the planner's view \citep{taxbreak2026,memoryfloor2026}.
KV tensors get a separate row: lifetime annotation \texttt{span=context} rather than \texttt{span=step} \citep{kwon2023vllm,dao2023flashdecoding}.

\paragraph{Step 2 ranking (how-to).}
Compute reuse count as number of on-chip loads per decode step (not per training step) \citep{patel2025llminference,dao2022flashattention}.
Divide by byte size; sort descending \citep{dlbook}.
Ties break toward tensors with longer live intervals---they exert sustained pressure \citep{wu2024mirage,yirage2026}.
The sorted list is the allocation priority queue for step~4 \citep{rotem2018glow,chen2020compilerbasedhardw}.

\paragraph{Step 3 peak walk (how-to).}
Simulate program order: on tensor birth add bytes to live set; on last use subtract \citep{yirage2026,rotem2018glow}.
When ping-pong applies, treat both buffers as live during the overlap window \citep{semianalysis2024tma,dao2022flashattention}.
Track maximum live-set size; that maximum is the number the memory pass compares to tier caps \citep{nvidia2022hopper,amd2024aie,chen2020compilerbasedhardw}.
If the walk is hand-done, cross-check with compiler pass output---discrepancies mean missing lifetimes \citep{yirage2026,wu2024mirage}.

\paragraph{Lifetime timeline discipline.}
Draw intervals on a horizontal time axis for one decode step---overlapping intervals reveal peak pressure visually \citep{wu2024mirage,dlbook}.
Reviewers should reject plans that show three large tensors overlapping without a documented split \citep{yirage2026,taxbreak2026}.

\begin{figure}[t]
\centering
\begin{tikzpicture}[vis base, scale=0.9, transform shape]
 \draw[vis arrow] (-4.5,0) -- (4.5,0);
 \foreach \x/\lab in {-4/0,-2/1,0/2,2/3,4/4} {
  \node[font=\scriptsize, visText] at (\x,-0.35) {$t_\lab$};
 }
 \draw[thick, blue!70] (-3.8,0.5) -- (-0.2,0.5) node[midway, above, font=\scriptsize] {QKV tile};
 \draw[thick, orange!80] (-2.5,1.0) -- (1.2,1.0) node[midway, above, font=\scriptsize] {softmax state};
 \draw[thick, green!60!black] (-1.0,1.5) -- (3.5,1.5) node[midway, above, font=\scriptsize] {MLP scratch};
 \draw[dashed, red!70] (0.8,-0.7) -- (0.8,2.0) node[below, font=\scriptsize, visText] {peak};
\end{tikzpicture}
\caption{Lifetime overlap within one decode step: peak footprint occurs where intervals stack.}
\label{fig:ch06_lifetime_timeline}
\end{figure}

\paragraph{Numeric walkthrough (Llama-3.2-1B layer, batch-1).}
Hidden size $d{=}2048$ in BF16: each activation vector is $4d{=}8$\,KB \citep{taxbreak2026,kwon2023vllm}.
Suppose a fused region keeps one hidden vector, one QKV weight tile (${\sim}32$\,KB for a thin tile), softmax carriers (${\sim}2$\,KB), and one MLP intermediate (${\sim}8$\,KB) simultaneously live during the attention peak \citep{dao2022flashattention}.
Single-buffer sum ${\approx}50$\,KB; with QKV ping-pong, add another $32$\,KB $\rightarrow$ peak ${\approx}82$\,KB before KV stream tiles \citep{semianalysis2024tma}.
Compare that peak to the GPU shared-memory budget from Chapter~\ref{chap:ch04}---legal with headroom for layout padding \citep{nvidia2022hopper,yirage2026}.
On XDNA, the same peak must fit tile SRAM caps from Chapter~\ref{chap:ch05}; if not, shrink tile width or split at the MLP export edge \citep{amd2024xdna,amd2024aie}.
This walkthrough is the spreadsheet behind Table~\ref{tab:ch06_peak_footprint}---not a substitute for your model's dimensions \citep{yirage2026,memoryfloor2026}.

\paragraph{Glow and Mirage parallels.}
Static buffer reuse across non-overlapping lifetimes is how classic compilers eliminate copies on low-level IR \citep{rotem2018glow}.
Mirage-style search proposes aggressive fusion orders; residency planning certifies which orders are \emph{legal} on each chip model \citep{wu2024mirage,yirage2026}.
Treat peak footprint as an invariant search rejects---not a post-run surprise \citep{mpk2025,chen2020compilerbasedhardw}.

\paragraph{Chapter~2 export-edge link.}
Iron Rule~2 (keep activations on-chip across stages) is implemented as short lifetimes between fused ops \citep{chen2020compilerbasedhardw,memoryfloor2026}.
Every export edge is a lifetime extension plus a bytes/token bill \citep{taxbreak2026,dao2022flashattention}.

\paragraph{$\checkmark$ Practice checkpoint.}
Attach a lifetime overlap sketch and peak-footprint worksheet before any fusion merge \citep{yirage2026,taxbreak2026}.

\section{Tri-Hardware Residency Adaptation}
\label{sec:ch06_tri_hardware_residency}

The four-step workflow is universal; \emph{tiers and control knobs} differ per class \citep{patel2025llminference,dlbook}.
This section states adaptation rules only---hardware deep dives live in Chapters~3--5 \citep{nvidia2022hopper,amd2024xdna,amd2024aie}.
Read it after completing a draft worksheet: you are mapping priorities to mechanisms, not re-learning silicon \citep{chen2020compilerbasedhardw,yirage2026}.

\paragraph{GPU in one paragraph.}
Programmer-managed shared memory and registers hold the fused decode working set; async copy engines overlap KV streams; L2 is contingency \citep{nvidia2022hopper,semianalysis2024tma,dao2022flashattention}.
Success metric: peak shared bytes plus layout fit Chapter~4 caps with no local-memory spill \citep{memoryfloor2026,yirage2026}.

\paragraph{CPU in one paragraph.}
Caches are implicit---you earn residency via blocking and NUMA placement \citep{dlbook,rotem2018glow,patel2025llminference}.
Success metric: weight and KV tiles reused inside L2/L3 before eviction; threads colocated with weight pages \citep{kwon2023vllm,memoryfloor2026}.

\paragraph{NPU in one paragraph.}
Tile and column SRAM are explicit; DMA schedules are the program \citep{amd2024xdna,amd2024aie}.
Success metric: static graph compiles with ping-pong included; DDR bytes/token matches roofline expectation \citep{patel2025llminference,taxbreak2026}.

\begin{figure}[t]
\centering
\begin{tikzpicture}[vis base, scale=0.82, transform shape]
 \node[font=\scriptsize\bfseries] at (-3.6,2.1) {GPU};
 \node[font=\scriptsize\bfseries] at (0,2.1) {CPU};
 \node[font=\scriptsize\bfseries] at (3.6,2.1) {NPU};
 \node[vis arch node, minimum width=2.0cm] at (-3.6,1.2) {\footnotesize Registers};
 \node[vis arch node, minimum width=2.0cm] at (-3.6,0.4) {\footnotesize Shared};
 \node[vis arch node, minimum width=2.0cm] at (-3.6,-0.4) {\footnotesize L2};
 \node[vis arch node, minimum width=2.0cm] at (0,1.2) {\footnotesize L1};
 \node[vis arch node, minimum width=2.0cm] at (0,0.4) {\footnotesize L2};
 \node[vis arch node, minimum width=2.0cm] at (0,-0.4) {\footnotesize L3};
 \node[vis arch node, minimum width=2.0cm] at (3.6,1.2) {\footnotesize Tile SRAM};
 \node[vis arch node, minimum width=2.0cm] at (3.6,0.4) {\footnotesize Column L2};
 \node[vis arch node, minimum width=2.0cm] at (3.6,-0.4) {\footnotesize DDR stream};
\end{tikzpicture}
\caption{Tri-hardware fast tiers (conceptual alignment; capacities in Ch.~3--5).}
\label{fig:ch06_tri_hw_tiers}
\end{figure}

\begin{table}[t]
\centering
\caption{Residency control model by hardware class (decode).}
\label{tab:ch06_residency_tiers}
\begin{tabular}{lll}
\toprule
\textbf{Class} & \textbf{Fast tier (budget input)} & \textbf{Who decides placement} \\
\midrule
GPU & Shared mem + registers (Ch.~4) & Programmer + layout \\
CPU & L1/L2/L3 via blocking (Ch.~3) & Access pattern + NUMA \\
NPU & Tile / column SRAM (Ch.~5) & Compiler DMA graph \\
\bottomrule
\end{tabular}
\end{table}

\subsection{GPU adaptation}
\label{sec:ch06_gpu_residency}

% AUTO_VISUAL:fig_gpu_residency_architecture
\begin{figure}[t]
\centering
\begin{tikzpicture}[vis base, scale=0.88, transform shape]
    \node[vis arch node] (s0) at (-5.03,0) {\footnotesize Host\\};
    \node[vis arch node] (s1) at (-1.68,0) {\footnotesize On-chip buffer\\};
    \node[vis arch node] (s2) at (1.67,0) {\footnotesize Compute array\\};
    \node[vis arch node] (s3) at (5.03,0) {\footnotesize Memory tier\\};
    \draw[vis arrow] (s0.east) -- (s1.west);
    \draw[vis arrow] (s1.east) -- (s2.west);
    \draw[vis arrow] (s2.east) -- (s3.west);
\end{tikzpicture}
\caption{GPU residency data path: programmer-managed on-chip buffers feed the compute array; overflow exports to slower tiers (capacities in Ch.~4).}
\label{fig:gpu_residency_architecture}
\end{figure}

\paragraph{Design focus (not hardware recap).}
Explicit shared-memory and register assignment; fusion region must fit SM budget from Chapter~\ref{chap:ch04} \citep{nvidia2022hopper,yirage2026}.
Bank-conflict-free layout is part of residency---see Chapter~\ref{chap:ch04} for swizzle rules \citep{semianalysis2024tma}.

\paragraph{GPU-specific methods.}
Trade occupancy for residency at batch-1 when parallelism cannot hide latency \citep{taxbreak2026,memoryfloor2026}.
Use async copy (TMA) to overlap KV streams---does not reduce bytes; fusion does \citep{memoryfloor2026,semianalysis2024tma}.
L2 is fallback, not plan---still aim to keep hot set in shared memory \citep{nvidia2022hopper}.

\paragraph{Applicable when.}
Fused MegaKernel with known peak $\le$ usable shared high-water (Chapter~4 byte sheet) \citep{yirage2026}.

\paragraph{Avoid when.}
Fusing past spill threshold---local memory traffic means residency plan failed \citep{memoryfloor2026,patel2025llminference}.

\paragraph{Register versus shared split.}
Place $(m,\ell,o)$ softmax carriers and thin accumulators in registers when register budget allows \citep{dao2022flashattention,nvidia2022hopper}.
Stage larger tiles (QKV, MLP fragments) in shared memory with conflict-free layout---Chapter~\ref{chap:ch04} covers bank swizzling \citep{semianalysis2024tma,yirage2026}.
When register pressure threatens spill, move carriers to shared before dropping fusion depth \citep{memoryfloor2026,taxbreak2026}.

\paragraph{Worked occupancy note.}
At batch-1, lowering occupancy to keep fusion in shared memory is often correct---parallelism cannot hide latency anyway \citep{taxbreak2026,memoryfloor2026,patel2025llminference}.
Document occupancy choice against peak footprint worksheet, not kernel microbench alone \citep{yirage2026,nvidia2022hopper}.

\paragraph{Layout is part of residency.}
A tile placed in shared memory but strided across the same bank serializes into multiple transactions---measurable bytes/token at batch-1 \citep{semianalysis2024tma,nvidia2022hopper}.
Tile shapes and padding are chosen with the allocation, not patched after profiling \citep{yirage2026,dao2022flashattention}.
Chapter~\ref{chap:ch04} documents swizzle rules; this chapter requires you to \emph{allocate} with them \citep{semianalysis2024tma}.

\paragraph{FlashAttention as residency pattern, not magic.}
Flash-style attention is the canonical example of keeping query tiles, KV tiles, and softmax statistics on chip instead of materializing scores to HBM \citep{dao2022flashattention}.
The lesson for decode is structural: identify the tensors with highest reuse-per-byte and pin them in the fastest legal tier for their lifetime \citep{dlbook,patel2025llminference}.
TMA overlap hides KV stream latency; it does not reduce bytes---fusion and lifetime shortening do \citep{memoryfloor2026,semianalysis2024tma}.

\subsection{CPU adaptation}
\label{sec:ch06_cpu_residency}

\paragraph{Design focus.}
Cache is hardware-managed---residency via blocking so working set reuses before eviction \citep{dlbook,rotem2018glow}.
\textbf{NUMA} affinity: pin decode threads to the node owning weights; remote memory can halve effective bandwidth on memory-bound loops \citep{patel2025llminference,memoryfloor2026}.

\paragraph{CPU-specific methods.}
Block GEMV/attention tiles to L1/L2; replicate read-only weights per node when cheaper than remote reads \citep{kwon2023vllm,dlbook}.
Validate with cache-miss counters at decode operating point---static capacity alone is insufficient \citep{patel2025llminference}.

\paragraph{Debug.}
If bytes/token is fine but ms/token regresses, check NUMA placement and prefetch pollution evicting hot tiles \citep{dlbook,taxbreak2026}.

\paragraph{Blocking recipe (decode).}
Choose block size $B$ so weight tile $+$ KV slice $+$ activations $\lesssim$ target L2 budget \citep{dlbook,kwon2023vllm}.
Reuse the block across the inner dimension before advancing to the next KV position \citep{dao2022flashattention,patel2025llminference}.
Residency here is \emph{algorithmic}---you do not name L2 addresses, you shape access \citep{rotem2018glow,yirage2026}.

\paragraph{Bandwidth arithmetic (why NUMA matters).}
At batch-1 a dense model re-reads essentially the full weight set every token, so token rate is bounded by weight bytes divided by effective memory bandwidth \citep{kwon2023vllm,memoryfloor2026}.
A NUMA-oblivious schedule---weights on node~0, threads on node~1---can halve effective bandwidth with flat cache-miss counters \citep{patel2025llminference,dlbook}.
The fix is placement and replication policy, not a faster inner loop \citep{taxbreak2026,memoryfloor2026}.

\paragraph{Prefetchers and empirical validation.}
Hardware prefetchers speculatively pull streams into cache---helpful for weight scans, harmful if they evict a small hot tensor you intended to keep resident \citep{dlbook}.
CPU residency tuning therefore pairs static blocking with miss-rate measurement at the decode operating point \citep{patel2025llminference,yirage2026}.
Contrast with GPU: you name shared addresses explicitly; on CPU you \emph{earn} residency through access shape \citep{rotem2018glow,chen2020compilerbasedhardw}.

\subsection{NPU adaptation}
\label{sec:ch06_npu_residency}

\paragraph{Design focus.}
No cache fallback---only compiler-placed SRAM plus DMA schedule (Chapter~\ref{chap:ch05}) \citep{amd2024xdna,amd2024aie}.
Illegal peak footprint fails compile or DDR-binds silently \citep{yirage2026,memoryfloor2026}.

\paragraph{NPU-specific methods.}
Size tiles to fit tile SRAM with ping-pong accounted; use column SRAM and spatial pipelines before DDR export \citep{amd2024aie,dao2022flashattention}.
Shape buckets and static graphs are residency prerequisites---not optional packaging \citep{patel2025llminference,amd2024xdna}.

\paragraph{Debug.}
DDR bytes/token rising with flat MAC active \% signals residency failure, not ``slow NPU'' \citep{memoryfloor2026,taxbreak2026}.

\paragraph{Column versus tile SRAM.}
Hot per-step tensors target tile SRAM; longer-lived KV slices may stream through column L2 with DMA overlap \citep{amd2024aie,amd2024xdna,dao2023flashdecoding}.
Misplacing a cold buffer in tile SRAM evicts hot tensors---anti-pattern~5 in Section~\ref{sec:ch06_residency_pitfalls} \citep{yirage2026,patel2025llminference}.

\paragraph{Double buffering without cache safety net.}
On spatial NPUs there is no hardware cache to absorb a late DMA load---compute tiles idle until operands arrive \citep{amd2024aie,amd2024xdna}.
Ping-pong between two tile buffers is therefore mandatory for streamed operands, and it halves usable capacity for those operands \citep{patel2025llminference}.
Step~3 of the workflow must count $2\times$ tile bytes \emph{before} codegen; Chapter~\ref{chap:ch05} lists the caps \citep{yirage2026,amd2024xdna}.

\paragraph{Spatial pipelines across tiles.}
Neighboring compute tiles can pass data through cascade links without DDR round-trips---a residency dimension GPUs approximate with shared memory and NPUs expose explicitly \citep{amd2024aie,dao2022flashattention}.
A strong NPU plan maps fused attention as a spatial sweep: inputs and final outputs touch DDR; intermediate scores stay in on-chip SRAM \citep{memoryfloor2026,amd2024xdna}.

\paragraph{Cross-class lesson.}
Identical fused math wants three different schedules: GPU explicit tiers, CPU blocking+NUMA, NPU static DMA \citep{chen2020compilerbasedhardw,yirage2026,patel2025llminference}.
A green GPU build proves nothing about CPU or NPU residency \citep{amd2024aie,memoryfloor2026,taxbreak2026}.

\section{Cross-Hardware Worked Example: Llama-3.2-1B Layer}
\label{sec:ch06_cross_hw_worked_example}

Setup: batch-1, ctx${=}2048$, hidden $d{=}2048$, one decoder layer, same high-level IR export \citep{taxbreak2026,kwon2023vllm,yirage2026}.
This is the same model family used in Chapter~4 and Chapter~5 worked cells---here we show \emph{allocation} rather than end-to-end goodput \citep{memoryfloor2026,taxbreak2026}.
The IR lists tensors; residency metadata assigns each tensor a tier per backend \citep{yirage2026,chen2020compilerbasedhardw}.

\begin{table}[t]
\centering
\small
\caption{Same tensors, different residency tiers (Llama-3.2-1B decode layer).}
\label{tab:ch06_tri_hw_allocation}
\begin{tabular}{@{}llll@{}}
\toprule
\textbf{Tensor} & \textbf{GPU (H100)} & \textbf{CPU} & \textbf{NPU (XDNA)} \\
\midrule
Input hidden state & Shared mem & L2 cache & Tile SRAM \\
QKV weight tile & Shared (TMA ping-pong) & L1 blocked tile & Tile SRAM (DMA ping-pong) \\
Softmax carriers $(m,\ell,o)$ & Registers & Registers & Tile SRAM slots \\
KV cache slice & Shared stream tile & L3 blocked stream & Column L2 stream \\
MLP intermediate & Shared mem & L2 cache & Tile SRAM \\
\bottomrule
\end{tabular}
\end{table}

\begin{table}[t]
\centering
\small
\caption{Peak footprint and fusion legality (illustrative).}
\label{tab:ch06_peak_footprint}
\begin{tabular}{@{}lccc@{}}
\toprule
\textbf{Metric} & \textbf{GPU} & \textbf{CPU} & \textbf{NPU} \\
\midrule
Peak fast-tier bytes & ${\sim}168$\,KB shared & ${\sim}2$\,MB L2 working set & ${\sim}34$\,KB tile \\
Double-buffer included? & Yes (QKV) & Yes (weights) & Yes (K/V) \\
Single fused layer legal? & Yes & Yes (blocked) & Yes (static graph) \\
Split required? & No & No & No at this $d$ \\
\bottomrule
\end{tabular}
\end{table}

\paragraph{Reading the tables.}
GPU places carriers in registers; NPU requires fixed carrier \emph{slots} in tile SRAM---same numerics, different legality mechanics \citep{dao2022flashattention,amd2024aie,nvidia2022hopper}.
CPU does not ``allocate L2''---blocking makes reuse likely; measure to verify \citep{dlbook,patel2025llminference}.
Triplet CI should compile all three or document which tier overflow forced a split \citep{yirage2026,chen2020compilerbasedhardw}.

\paragraph{Worked contrast.}
Chapter~4 Table~\ref{tab:ch04_worked_results} and Chapter~5 Table~\ref{tab:ch05_worked_results} report goodput cells for single hardware \citep{memoryfloor2026,taxbreak2026}.
This table explains \emph{why} those cells differ---residency plans diverge while IR stays shared \citep{yirage2026,patel2025llminference}.

\paragraph{Split-fusion trigger example.}
Raise hidden size or head count until GPU peak shared memory exceeds Chapter~4 cap---memory pass should split at MLP export edge \citep{yirage2026,nvidia2022hopper,memoryfloor2026}.
Same IR on NPU may split earlier due to smaller tile SRAM---triplet tests expose the gap \citep{amd2024aie,amd2024xdna,chen2020compilerbasedhardw}.

\paragraph{Ablation discipline.}
Document unfused, partially fused, and fully fused peak footprints per target---label which tier overflow forced each split \citep{wu2024mirage,yirage2026,taxbreak2026,memoryfloor2026}.
Reviewers reject ms/token wins that hide demoted tensors \citep{patel2025llminference}.

\paragraph{Step-by-step allocation narrative.}
\textbf{GPU:} assign softmax carriers to registers; stage QKV and hidden vectors in shared memory with ping-pong on the weight tile; stream KV slices through TMA into the same fusion region \citep{dao2022flashattention,semianalysis2024tma,nvidia2022hopper}.
Run step~3 peak walk---if shared high-water exceeds Chapter~4 usable cap, split at MLP export \citep{yirage2026,memoryfloor2026}.
\textbf{CPU:} pick block size so weight tile $+$ KV slice fits L2; pin threads to the NUMA node that owns weights; verify with LLC miss counters \citep{dlbook,patel2025llminference}.
\textbf{NPU:} lower QKV tile width until tile SRAM peak (with ping-pong) fits Chapter~5 budget; schedule column-L2 streaming for KV history \citep{amd2024xdna,amd2024aie}.
Archive three worksheets---identical IR, divergent peaks \citep{yirage2026,chen2020compilerbasedhardw}.

\paragraph{Connection to serving stacks.}
vLLM-style paging solves \emph{where KV lives in DRAM}; this chapter solves \emph{which KV bytes enter the fast tier each step} \citep{kwon2023vllm,dao2023flashdecoding}.
Both layers matter: paging without on-chip residency still pays HBM round-trips per attention op \citep{taxbreak2026,memoryfloor2026}.

\begin{table}[t]
\centering
\small
\caption{Peak footprint line items (GPU, Llama-3.2-1B layer, illustrative).}
\label{tab:ch06_gpu_peak_line_items}
\begin{tabular}{@{}lrr@{}}
\toprule
\textbf{Buffer} & \textbf{Bytes} & \textbf{Ping-pong?} \\
\midrule
Hidden activation (BF16) & 8\,KB & No \\
QKV weight tile & 32\,KB & Yes ($\times 2$) \\
Softmax carriers & 2\,KB & No \\
KV stream tile & 16\,KB & Yes ($\times 2$) \\
MLP intermediate & 8\,KB & No \\
\midrule
\textbf{Peak sum} & {\footnotesize $\sim$168}\,KB & \\
\bottomrule
\end{tabular}
\end{table}

\paragraph{Line-item discipline.}
Table~\ref{tab:ch06_gpu_peak_line_items} is how Pattern~1 worksheets should look---every row auditable \citep{yirage2026,chen2020compilerbasedhardw}.
Summing without rows hides double-buffer mistakes \citep{semianalysis2024tma,memoryfloor2026}.
Cross-check the total against Chapter~4 shared-memory budget before claiming fusion legality \citep{nvidia2022hopper,taxbreak2026}.

\paragraph{$\checkmark$ Practice checkpoint.}
Fill Table~\ref{tab:ch06_tri_hw_allocation} for your model before claiming triplet parity \citep{yirage2026,amd2024xdna,nvidia2022hopper}.

\section{YiRage Memory Pass: Residency Rules as Compile Gates}
\label{sec:ch06_yirage_memory_pass}

\paragraph{Advance notice.}
Compiler/IR literacy helps here; kernel tuners should skim invariants and return for pass I/O when debug fails \citep{chen2020compilerbasedhardw,yirage2026}.

\textbf{YiRage} encodes per-class residency contracts as IR metadata, then verifies schedules before codegen \citep{yirage2026}.
\textbf{Pipeline position:} after bufferization, before target codegen---legality pass, not optional lint \citep{rotem2018glow,wu2024mirage}.

\begin{figure}[t]
\centering
\begin{tikzpicture}[vis base, scale=0.82, transform shape]
 \node[vis pipeline node] (buf) at (-3.2,0) {\footnotesize Bufferize};
 \node[vis arch node] (mem) at (0,0) {\footnotesize Memory\\pass};
 \node[vis pipeline node] (cg) at (3.2,0) {\footnotesize Codegen};
 \draw[vis arrow] (buf) -- (mem) -- (cg);
 \node[font=\scriptsize, visText] at (0,-1.0) {Peak footprint vs tier cap; split or demote};
\end{tikzpicture}
\caption{YiRage memory pass in the lowering pipeline.}
\label{fig:ch06_yirage_memory_pass}
\end{figure}

\begin{table}[t]
\centering
\small
\caption{YiRage memory pass: inputs, checks, outputs.}
\label{tab:ch06_memory_pass_io}
\begin{tabular}{@{}p{2.2cm}p{4.2cm}p{3.6cm}@{}}
\toprule
\textbf{Dimension} & \textbf{Check} & \textbf{On failure} \\
\midrule
Peak footprint & Simultaneous live buffers + 2$\times$ ping-pong & Split fusion or demote tensor \\
GPU layout & Bank conflict hazard (Ch.~4 rules) & Re-tile or pad \\
CPU NUMA & Weight/KV node affinity metadata & Emit warning or reject \\
NPU DMA & Column bandwidth vs pipeline depth & Reschedule DMA chain \\
Superoptimizer & Mirage proposal vs tier cap & Reject illegal fusion \citep{wu2024mirage} \\
\bottomrule
\end{tabular}
\end{table}

\paragraph{Input/output contract.}
\textbf{Input:} memref IR with lifetime intervals and target chip model from Chapter~3 \citep{chen2020compilerbasedhardw,yirage2026}.
\textbf{Output:} legal tier assignment per buffer \emph{or} structured error citing overflow row and suggested split edge \citep{amd2024aie,nvidia2022hopper,patel2025llminference}.
Metadata fields typically include: \texttt{lifetime\_start}, \texttt{lifetime\_end}, \texttt{reuse\_score}, \texttt{preferred\_tier}, \texttt{ping\_pong}, and \texttt{demotion\_cost\_bytes} \citep{yirage2026,rotem2018glow}.
Downstream codegen reads tier assignment; it does not re-infer lifetimes \citep{wu2024mirage,chen2020compilerbasedhardw}.

\paragraph{Interaction with search.}
Mirage-style superoptimizers propose aggressive fusions; memory pass certifies residency before emit \citep{wu2024mirage,mpk2025,yirage2026}.
Search explores performance; pass enforces mechanical sympathy---separation that makes heterogeneous fleets deployable \citep{patel2025llminference,memoryfloor2026}.

\paragraph{Error message contract.}
Overflow reports must cite: tensor name, live interval, bytes, tier cap, and recommended split edge \citep{yirage2026,amd2024aie,nvidia2022hopper}.
Opaque ``compile failed'' messages recreate GPU-porting pain on NPU \citep{patel2025llminference,chen2020compilerbasedhardw}.

\paragraph{Regression tests.}
CI should fail when a new buffer appears in the hot path without lifetime metadata \citep{yirage2026,rotem2018glow}.
Golden tests compare peak footprint against stored worksheets for canonical layers \citep{wu2024mirage,mpk2025}.

\paragraph{Closed loop with superoptimization.}
A practical lowering loop alternates: Mirage proposes a fused schedule $\rightarrow$ memory pass computes peak footprint and tier assignment $\rightarrow$ illegal proposals are rejected with a split hint $\rightarrow$ search retries with a narrower fusion region \citep{wu2024mirage,yirage2026,mpk2025}.
This separation keeps search ambitious and deployment safe---the pass is the mechanical-sympathy gate Chapter~3 promised \citep{chen2020compilerbasedhardw,patel2025llminference}.
Kernel engineers should treat pass errors as \emph{design feedback}, not annoyances to bypass with smaller tiles alone \citep{memoryfloor2026,taxbreak2026}.

\paragraph{Overflow strategies (two sanctioned paths).}
When peak exceeds tier cap, YiRage may (a)~split fusion at the lowest reuse-per-byte export edge, preserving carriers on chip, or (b)~demote a low-score tensor to the next tier and attach an expected bytes/token cost to the IR \citep{yirage2026,chen2020compilerbasedhardw}.
Undocumented demotion is how teams ship ``fast'' kernels that regress on longer contexts \citep{patel2025llminference,memoryfloor2026}.
Pick (a) when split restores legality with acceptable launch increase; pick (b) only when demotion cost is below serving SLO \citep{taxbreak2026,wu2024mirage}.
Both paths must emit human-readable rationale for review packets \citep{yirage2026,chen2020compilerbasedhardw}.

\paragraph{$\checkmark$ Practice checkpoint.}
Attach memory-pass legality report (or overflow row) on every fusion PR \citep{yirage2026,chen2020compilerbasedhardw}.

\section{Residency Engineering Patterns}
\label{sec:ch06_residency_patterns}

The four-step workflow becomes repeatable when encoded as review patterns \citep{chen2020compilerbasedhardw,yirage2026}.
Use these in design docs the same way Chapter~5 uses XDNA patterns---Scene, Solution, Benefit, Notes \citep{patel2025llminference,amd2024aie}.

\paragraph{Pattern 1: Peak worksheet before fusion merge.}
\textbf{Scene:} team proposes merging attention and MLP into one kernel \citep{taxbreak2026}.
\textbf{Solution:} run step~1--3 on the fused DAG; attach live-set timeline (Figure~\ref{fig:ch06_lifetime_timeline}) \citep{wu2024mirage,dlbook}.
\textbf{Benefit:} catches overflow before codegen and before NPU compile \citep{yirage2026,amd2024xdna}.
\textbf{Notes:} include ping-pong bytes; compare to Chapter~3 chip-model row \citep{chen2020compilerbasedhardw}.

\paragraph{Pattern 2: Carrier-first tier assignment.}
\textbf{Scene:} softmax state $(m,\ell,o)$ contends with QKV tiles for registers or tile SRAM \citep{dao2022flashattention}.
\textbf{Solution:} assign carriers to the fastest tier first; demote low-reuse scratch before touching carriers \citep{dao2022flashattention,patel2025llminference}.
\textbf{Benefit:} preserves numerically sensitive state on chip across long contexts \citep{dao2023flashdecoding}.
\textbf{Notes:} GPU $\rightarrow$ registers; NPU $\rightarrow$ fixed slots---see Table~\ref{tab:ch06_tri_hw_allocation} \citep{amd2024aie,nvidia2022hopper}.

\paragraph{Pattern 3: KV stream with explicit overlap budget.}
\textbf{Scene:} KV history too large to resident; each token loads a slice \citep{kwon2023vllm,dao2023flashdecoding}.
\textbf{Solution:} double-buffer KV tiles; schedule loads to complete before consumers stall \citep{semianalysis2024tma,amd2024xdna}.
\textbf{Benefit:} hides DDR latency without pretending KV ``fits'' in tile SRAM \citep{patel2025llminference,memoryfloor2026}.
\textbf{Notes:} step~3 must count $2\times$ KV tile bytes during overlap window \citep{yirage2026,dao2022flashattention}.

\paragraph{Pattern 4: Export-edge split discipline.}
\textbf{Scene:} fused region overflows fastest tier \citep{memoryfloor2026,yirage2026}.
\textbf{Solution:} split at Chapter~2 export edge with lowest reuse-per-byte tensor \citep{chen2020compilerbasedhardw,taxbreak2026}.
\textbf{Benefit:} preserves the highest-value fusion while restoring legality \citep{wu2024mirage,mpk2025}.
\textbf{Notes:} memory pass should emit the recommended edge in the error text \citep{yirage2026}.

\paragraph{Pattern 5: Triplet worksheet parity.}
\textbf{Scene:} CUDA schedule ``works''; Ryzen AI compile fails \citep{amd2024xdna,taxbreak2026}.
\textbf{Solution:} fill Table~\ref{tab:ch06_tri_hw_allocation} for all three backends before merge \citep{yirage2026,chen2020compilerbasedhardw}.
\textbf{Benefit:} prevents GPU-only residency assumptions in shared IR \citep{patel2025llminference,memoryfloor2026}.
\textbf{Notes:} peak bytes differ; IR parity $\neq$ residency parity \citep{amd2024aie,nvidia2022hopper}.

\paragraph{Pattern 6: Bytes/token proof artifact.}
\textbf{Scene:} ms/token improved but serving cost unclear \citep{memoryfloor2026,taxbreak2026}.
\textbf{Solution:} archive bytes/token worksheet alongside latency; note any demoted tensors \citep{patel2025llminference,yirage2026}.
\textbf{Benefit:} separates real residency wins from measurement noise \citep{dao2022flashattention,memoryfloor2026}.
\textbf{Notes:} tie to Chapter~1 metric definitions \citep{taxbreak2026}.

\paragraph{Pattern usage in review.}
Require Pattern~1 worksheet on fusion PRs; Pattern~5 triplet row on IR changes; Pattern~6 bytes proof on performance claims \citep{yirage2026,chen2020compilerbasedhardw,patel2025llminference}.

\section{Residency Design Review Workflow}
\label{sec:ch06_residency_review}

Production teams benefit from a fixed review sequence---same spirit as Chapter~5 review gates, applied to cross-hardware residency \citep{patel2025llminference,yirage2026,amd2024xdna}.

\paragraph{Step R1: IR freeze.}
Freeze the fused subgraph under review: list inputs, outputs, and export edges \citep{chen2020compilerbasedhardw,taxbreak2026}.
No residency review on a moving graph \citep{wu2024mirage,yirage2026}.

\paragraph{Step R2: Lifetime inventory.}
For each intermediate, record producer op, last consumer, and byte size \citep{rotem2018glow,dlbook}.
Flag KV tensors separately---lifetime spans context, not one step \citep{kwon2023vllm,dao2023flashdecoding}.

\paragraph{Step R3: Peak worksheets (triplet).}
Run step~3 for GPU, CPU, and NPU chip models from Chapter~3 \citep{chen2020compilerbasedhardw,nvidia2022hopper,amd2024aie}.
Attach three numbers or mark split/demotion explicitly \citep{yirage2026,memoryfloor2026}.

\paragraph{Step R4: Memory-pass dry run.}
Run YiRage memory pass on each backend; archive pass report or overflow row \citep{yirage2026,wu2024mirage}.
Reject merges with opaque compile failures \citep{patel2025llminference,amd2024xdna}.

\paragraph{Step R5: Checklist sign-off.}
Complete Section~\ref{sec:ch06_residency_pitfalls} checklist; link Pattern~1--6 artifacts \citep{chen2020compilerbasedhardw,taxbreak2026}.
Only then proceed to kernel tuning in Part~IV \citep{semianalysis2024tma,dao2022flashattention}.

\paragraph{Facilitation notes.}
Residency review is cross-functional: compiler authors own pass output, kernel engineers own tier justification, serving owners own bytes/token acceptance \citep{patel2025llminference,yirage2026,memoryfloor2026}.
Disputes resolve with peak worksheets, not opinions \citep{taxbreak2026,chen2020compilerbasedhardw}.

\paragraph{Artifact bundle template.}
A complete residency review packet contains: (1)~fused IR snippet with export edges highlighted, (2)~lifetime table (tensor, bytes, $t_{\mathrm{prod}}$, $t_{\mathrm{last}}$), (3)~three peak-footprint numbers, (4)~memory-pass log excerpt, (5)~signed checklist \citep{yirage2026,chen2020compilerbasedhardw,patel2025llminference}.
Store packets next to benchmark TSV rows so regressions bisect quickly \citep{taxbreak2026,memoryfloor2026}.
This is the operational form of ``residency is a contract''---not a slide deck \citep{dlbook,rotem2018glow}.

\section{Residency Pitfalls, Checklist, and Debug}
\label{sec:ch06_residency_pitfalls}

\subsection{Five common residency anti-patterns}
\begin{enumerate}
\item \textbf{Over-fusion.}
\emph{Symptom:} spill / local-memory loads on GPU; compile fail on NPU \citep{memoryfloor2026,yirage2026}.
\emph{Root cause:} fusion depth chosen for launch count, not peak footprint \citep{taxbreak2026,wu2024mirage}.
\emph{Fix:} split at Chapter~2 export edge; re-run peak footprint step~3 \citep{chen2020compilerbasedhardw}.
\emph{Prevention:} Pattern~1 worksheet mandatory before merge \citep{yirage2026}.

\item \textbf{Double-buffer not in budget.}
\emph{Symptom:} fits in spreadsheet, overflows after enabling ping-pong \citep{semianalysis2024tma,dao2022flashattention}.
\emph{Root cause:} step~3 counted single-buffer bytes only \citep{yirage2026,memoryfloor2026}.
\emph{Fix:} multiply streamed operand bytes by two in peak calculation \citep{yirage2026}.
\emph{Prevention:} mark \texttt{ping\_pong=true} in IR before pass \citep{rotem2018glow,chen2020compilerbasedhardw}.

\item \textbf{Stretched lifetimes.}
\emph{Symptom:} fast tier full despite small tensors \citep{wu2024mirage,dlbook}.
\emph{Root cause:} schedule interleaves unrelated ops between producer and consumer \citep{taxbreak2026,patel2025llminference}.
\emph{Fix:} reschedule fused DAG so high-reuse tensors have shorter live ranges \citep{taxbreak2026}.
\emph{Prevention:} co-optimize schedule with step~1 inventory \citep{wu2024mirage,yirage2026}.

\item \textbf{NUMA blindness (CPU).}
\emph{Symptom:} half effective bandwidth, flat cache miss rate \citep{patel2025llminference,memoryfloor2026}.
\emph{Root cause:} threads and weights on different sockets \citep{dlbook,kwon2023vllm}.
\emph{Fix:} pin threads; replicate weights per node \citep{dlbook}.
\emph{Prevention:} Gate~C in this section on every CPU benchmark \citep{patel2025llminference,taxbreak2026}.

\item \textbf{NPU tier misplacement.}
\emph{Symptom:} hot tensor in column L2 while cold tensor consumes tile SRAM \citep{amd2024aie,amd2024xdna}.
\emph{Root cause:} DMA schedule optimized for line count, not reuse score \citep{yirage2026,patel2025llminference}.
\emph{Fix:} reorder DMA schedule; demote cold buffers first \citep{yirage2026,patel2025llminference}.
\emph{Prevention:} Pattern~2 carrier-first assignment \citep{dao2022flashattention,amd2024aie}.
\end{enumerate}

\paragraph{Megakernel prerequisites (Part~IV).}
Chapters~7--11 assume residency is already legal \citep{taxbreak2026,mpk2025,wu2024mirage}.
\textbf{Chapter~7} static role assignment maps thread blocks to fused stages---roles are meaningless if the working set spills \citep{patel2025llminference,yirage2026}.
\textbf{Chapter~8} TMA choreography overlaps KV streams---ping-pong must already appear in step~3 worksheets \citep{semianalysis2024tma,dao2022flashattention}.
\textbf{Chapter~10} softmax carriers need fixed on-chip homes---Pattern~2 must have run \citep{dao2022flashattention,dao2023flashdecoding}.
\textbf{Chapter~11} full MegaKernel assembly is the capstone---only shippable when memory pass reports green on all fleet backends \citep{yirage2026,chen2020compilerbasedhardw,memoryfloor2026}.
Skipping this chapter turns Part~IV into microbenchmark tuning \citep{taxbreak2026,memoryfloor2026}.
If you are maintaining an existing CUDA kernel, use the tri-hardware table as a porting checklist before touching Ryzen AI backends or tuning MegaKernel schedules \citep{amd2024xdna,patel2025llminference,yirage2026,mpk2025}.

\subsection{Residency design self-check (review checklist)}
\begin{itemize}
\item[$\square$] All tensor lifetimes documented; peak footprint calculated
\item[$\square$] Double-buffer 2$\times$ occupancy included in peak
\item[$\square$] KV streaming overlap plan documented
\item[$\square$] Target fast-tier capacity not exceeded (Ch.~3--5 caps)
\item[$\square$] Demoted tensors have accepted bandwidth cost
\item[$\square$] CPU: NUMA affinity plan attached
\item[$\square$] NPU: DMA chain bandwidth sanity-checked
\item[$\square$] Triplet IR memory-pass report archived
\end{itemize}

\paragraph{Postmortem sequence.}
On residency regression: replay peak worksheet $\rightarrow$ diff bytes/token $\rightarrow$ check spill counters $\rightarrow$ bisect fusion split $\rightarrow$ re-run memory pass \citep{taxbreak2026,memoryfloor2026,yirage2026,patel2025llminference}.
Never tune MMA tiles before bytes/token replay \citep{nvidia2022hopper,amd2024aie}.

\subsection{Review gates (symptom $\rightarrow$ debug)}
\paragraph{Gate A: shared-memory overflow on GPU.}
\textbf{Symptom:} compile warning or runtime local-memory loads \citep{memoryfloor2026,nvidia2022hopper}.
\textbf{Debug:} re-run step~3 with ping-pong; check bank layout; compare to Chapter~4 cap \citep{semianalysis2024tma,yirage2026}.
\textbf{Fix:} split fusion or demote lowest-score scratch \citep{wu2024mirage,chen2020compilerbasedhardw}.

\paragraph{Gate B: NPU compile fail on tile SRAM.}
\textbf{Symptom:} DMA schedule cannot fit ping-pong tiles \citep{amd2024xdna,amd2024aie}.
\textbf{Debug:} print peak tile bytes from memory pass; verify column vs tile assignment \citep{yirage2026,patel2025llminference}.
\textbf{Fix:} shrink tile width or move KV stream to column L2 \citep{amd2024aie,dao2023flashdecoding}.

\paragraph{Gate C: CPU bandwidth cliff.}
\textbf{Symptom:} token rate half of roofline expectation; flat LLC misses \citep{patel2025llminference,memoryfloor2026}.
\textbf{Debug:} NUMA node of threads vs weights; remote access counters \citep{dlbook}.
\textbf{Fix:} pin threads; replicate weights per node \citep{kwon2023vllm,taxbreak2026}.

\paragraph{Gate D: silent demotion.}
\textbf{Symptom:} latency OK at ctx${=}512$, regresses at ctx${=}8192$ \citep{patel2025llminference,dao2023flashdecoding}.
\textbf{Debug:} diff memory-pass tier assignments across context buckets \citep{yirage2026,amd2024xdna}.
\textbf{Fix:} document demotion bytes cost; split fusion instead of hiding spill \citep{memoryfloor2026,chen2020compilerbasedhardw}.

\paragraph{Documentation contract.}
Every residency promotion ships: lifetime diagram, peak worksheet, tri-hardware table row, memory-pass report, checklist tick marks \citep{yirage2026,chen2020compilerbasedhardw,patel2025llminference,taxbreak2026,memoryfloor2026}.
Private spreadsheets do not scale across fleets \citep{patel2025llminference}.

\paragraph{Reader exercises.}
(1) Compute peak footprint for one layer with ping-pong. (2) Explain one anti-pattern symptom from your last merge. (3) Fill checklist for your model \citep{yirage2026,dao2022flashattention,amd2024xdna,nvidia2022hopper,taxbreak2026,memoryfloor2026}.

\paragraph{$\checkmark$ Practice checkpoint.}
Merge blocked until checklist complete or deferral documented with overflow row \citep{yirage2026,chen2020compilerbasedhardw}.

\section{Key Takeaways}
\label{sec:ch06_key_takeaways}

Residency planning is the hinge between hardware chapters and MegaKernel implementation \citep{patel2025llminference,yirage2026,chen2020compilerbasedhardw}.
The principles below are intentionally actionable---each pairs a design rule with a concrete next step your team can adopt in the next sprint \citep{taxbreak2026,memoryfloor2026}.

\begin{itemize}
\item \textbf{Plan-before-codegen principle.}
Residency is a contract between fused IR and silicon, not a post-hoc tweak after kernels compile \citep{dlbook,yirage2026}.
If peak footprint is unknown at merge time, you do not yet have a decode optimization---you have a prototype \citep{memoryfloor2026,taxbreak2026}.
\textbf{Action:} run the four-step workflow and attach worksheets before any codegen PR \citep{chen2020compilerbasedhardw,rotem2018glow}.

\item \textbf{Lifetime-first principle.}
Tier placement follows lifetime span and reuse-per-byte, not operator names or framework defaults \citep{dao2022flashattention,wu2024mirage}.
Shortening a high-reuse tensor's lifetime via scheduling often beats enlarging shared memory or tile SRAM \citep{taxbreak2026,patel2025llminference}.
\textbf{Action:} draw Figure~\ref{fig:ch06_lifetime_timeline} for each fused region; reject schedules that stack three large tensors without justification \citep{yirage2026,dlbook}.

\item \textbf{Hardware-adaptation principle.}
There is no portable residency template: GPU uses explicit shared memory, CPU uses blocking and NUMA, NPU uses static DMA graphs \citep{patel2025llminference,amd2024xdna,nvidia2022hopper}.
Identical IR with identical ms/token on one backend proves nothing about another \citep{yirage2026,chen2020compilerbasedhardw,memoryfloor2026}.
\textbf{Action:} fill Table~\ref{tab:ch06_tri_hw_allocation} for every backend in your fleet before claiming parity \citep{amd2024aie,taxbreak2026}.

\item \textbf{Double-buffer trade principle.}
Ping-pong buys pipeline overlap at $2\times$ space for streamed operands---worth it only when step~3 still fits \citep{semianalysis2024tma,dao2022flashattention}.
Silent omission of doubled bytes is anti-pattern~2 \citep{yirage2026,memoryfloor2026}.
\textbf{Action:} show doubled bytes explicitly in peak worksheets and IR \texttt{ping\_pong} flags \citep{rotem2018glow,chen2020compilerbasedhardw}.

\item \textbf{Verifiable-plan principle.}
Residency must be quantified, compiler-checked, and archived---experience alone does not scale across triplet fleets \citep{yirage2026,wu2024mirage,patel2025llminference}.
The memory pass is the mechanical gate; the checklist is the human gate \citep{chen2020compilerbasedhardw,mpk2025}.
\textbf{Action:} attach memory-pass reports and Section~\ref{sec:ch06_residency_pitfalls} checklist ticks to CI artifacts \citep{yirage2026,taxbreak2026}.
\end{itemize}

\paragraph{How this chapter connects.}
Chapter~3 gave constraint matrices; Chapters~4--5 gave per-class caps and mechanisms \citep{chen2020compilerbasedhardw,nvidia2022hopper,amd2024xdna}.
This chapter gave the procedure to turn caps into assignments \citep{yirage2026,dlbook}.
Chapter~\ref{chap:ch07} and beyond implement the schedules those assignments enable \citep{taxbreak2026,patel2025llminference,mpk2025}.

\section{Conclusion}
\label{sec:ch06_conclusion}

On-chip residency is where decode optimization becomes concrete: lifetimes, peak footprint, and per-class tier assignment \citep{dlbook,chen2020compilerbasedhardw}.
Chapters~4--5 gave hardware-specific bounds; this chapter gave the \emph{method} to use them without duplicating catalogs \citep{nvidia2022hopper,amd2024xdna,yirage2026}.
You now have a workflow, patterns, gates, and artifacts that turn those bounds into shippable fused kernels across a heterogeneous fleet \citep{patel2025llminference,taxbreak2026,memoryfloor2026}.

\paragraph{Closing metaphor.}
Think of residency as airport gate assignment: flights (tensors) differ in size and connection time (lifetime); gates (tiers) are finite \citep{dlbook}.
A schedule that minimizes runway taxiing (launches) but parks wide-body jets at regional gates (shared-memory overflow) still misses departure \citep{taxbreak2026,yirage2026}.
The four-step workflow is the operations desk; the memory pass is the automated overbooking alarm \citep{wu2024mirage,chen2020compilerbasedhardw}.

Chapter~\ref{chap:ch07} builds the next MegaKernel technique on this foundation: \textbf{static role assignment}---fixing each thread block's job at compile time so the resident working set stops moving \citep{taxbreak2026,patel2025llminference}.
Static roles only help when peak footprint is already legal---otherwise roles shuffle spill traffic \citep{memoryfloor2026,yirage2026}.
Carry peak-footprint worksheets and Table~\ref{tab:ch06_tri_hw_allocation} into every subsequent kernel chapter \citep{yirage2026,memoryfloor2026}.

\paragraph{Recap for practitioners.}
If you remember nothing else: decode optimization is residency optimization; lifetimes drive tiers; ping-pong doubles bytes; triplet fleets need triplet worksheets; the memory pass is the gate \citep{dao2022flashattention,patel2025llminference,yirage2026,taxbreak2026,memoryfloor2026}.
Run the four-step workflow once per fused region per backend before writing a single line of device code \citep{chen2020compilerbasedhardw,rotem2018glow}.
When in doubt, shorten lifetimes via scheduling before buying larger chips \citep{wu2024mirage,taxbreak2026}.

\paragraph{Part~IV bridge.}
Residency planning is prerequisite for TMA (Ch.~8), sync (Ch.~9), softmax carriers (Ch.~10), and full MegaKernel assembly (Ch.~11) \citep{dao2022flashattention,semianalysis2024tma,yirage2026}.
If peak footprint is unproved here, later chapters optimize the wrong problem \citep{taxbreak2026,memoryfloor2026}.

\paragraph{Summary before Part~IV.}
Lifetime method, tri-hardware adaptation, worked allocation, memory pass, engineering patterns, and pitfalls form the residency spine of the book \citep{dlbook,chen2020compilerbasedhardw,yirage2026,dao2022flashattention,patel2025llminference,amd2024xdna,nvidia2022hopper,taxbreak2026,memoryfloor2026,wu2024mirage,rotem2018glow,mpk2025}.
Bytes/token proof comes from planning---not from hoping caches behave \citep{memoryfloor2026,taxbreak2026}.

\paragraph{What to carry forward.}
Every subsequent kernel chapter assumes you can (1)~name each tensor's lifetime, (2)~compute peak footprint with ping-pong, (3)~justify tier assignment per backend, and (4)~attach a memory-pass or checklist artifact \citep{yirage2026,chen2020compilerbasedhardw,patel2025llminference}.
MegaKernel techniques in Part~IV add synchronization, TMA choreography, and softmax carrier placement---but they do not replace residency planning \citep{semianalysis2024tma,dao2022flashattention,taxbreak2026}.
If you skip the worksheets here, you will optimize launch count while bytes/token flatlines \citep{memoryfloor2026,taxbreak2026}.

\paragraph{Teaching versus production.}
In teaching mode, hand-draw lifetime timelines and compare tri-hardware tables on a whiteboard \citep{dlbook,wu2024mirage}.
In production mode, automate Pattern~1--6 in CI and treat memory-pass failures as merge blockers \citep{yirage2026,patel2025llminference,mpk2025}.
The book shows both: concepts in Sections~\ref{sec:ch06_lifetime_planning}--\ref{sec:ch06_cross_hw_worked_example}, gates in Sections~\ref{sec:ch06_yirage_memory_pass}--\ref{sec:ch06_residency_pitfalls} \citep{chen2020compilerbasedhardw,memoryfloor2026}.

\paragraph{Done criteria.}
You are done when you can compute peak footprint with ping-pong, explain tier differences for one tensor across GPU/CPU/NPU, and pass the residency checklist \citep{yirage2026,patel2025llminference,amd2024aie,nvidia2022hopper}.
Proceed to Chapter~\ref{chap:ch07} with worksheets attached \citep{chen2020compilerbasedhardw}.
Bring Table~\ref{tab:ch06_tri_hw_allocation}, Table~\ref{tab:ch06_gpu_peak_line_items}, and your signed checklist---Chapter~7 will assume those numbers are already consensus on your team \citep{yirage2026,patel2025llminference,taxbreak2026}.

\paragraph{Hands-on recap.}
(1) Run four-step workflow on one decoder layer. (2) Fill tri-hardware table and peak worksheet. (3) Run memory pass on each backend. (4) Complete checklist. (5) Archive bytes/token proof alongside latency \citep{taxbreak2026,memoryfloor2026,yirage2026,patel2025llminference,dao2022flashattention,amd2024xdna,nvidia2022hopper,semianalysis2024tma,dlbook,wu2024mirage,rotem2018glow,mpk2025,kwon2023vllm,dao2023flashdecoding,liu2024deepseekv3}.
If any step is missing, downstream MegaKernel work is premature \citep{chen2020compilerbasedhardw}.
Teams that skip worksheets often rediscover the same spill bugs in Chapters~8--11 under higher time pressure \citep{taxbreak2026,memoryfloor2026,semianalysis2024tma}.
Treat this recap as the exit criteria for Part~II residency literacy before entering Part~IV implementation depth \citep{patel2025llminference,yirage2026,mpk2025}.
Archive worksheets in the same repository as benchmark TSV files so bisects correlate code, residency, and bytes/token in one timeline \citep{taxbreak2026,memoryfloor2026,chen2020compilerbasedhardw}.
When onboarding a new backend, start from Pattern~5: fill the tri-hardware row before tuning kernels---porting pain is usually residency mismatch, not missing intrinsics \citep{yirage2026,amd2024xdna,nvidia2022hopper,patel2025llminference}.
Re-read Section~\ref{sec:industrial_hook} before your next fusion review---it states the thesis this chapter proves in worksheets and pass logs \citep{taxbreak2026,memoryfloor2026,dao2022flashattention}.
Part~II ends here: you know the silicon (Ch3--5) and how to plan bytes on it (Ch6) \citep{chen2020compilerbasedhardw,yirage2026,patel2025llminference}.
Enter Part~IV with worksheets, not hope, on every backend you ship in production \citep{memoryfloor2026,taxbreak2026,patel2025llminference}.

\typeout{END_CHAPTER "ch06" \theabspage}
